% !TEX root = main.tex
\chapter{Teoría Global de Superficies en $\bb{R}^3$}

\section{Introducción}

Recordamos algunas definiciones que ya se vieron en la asignatura de Geometría Global de Curvas y Superficies.

\begin{definicion}
    Sea $S\subset \bb{R}^3$ no vacío, diremos que $S$ es una \textbf{superficie} si para todo $p\in S$ existen
    \begin{enumerate}[label=\alph*)]
        \item \label{def:sup_exists1} $U\subset \bb{R}^2$ abierto.
        \item \label{def:sup_exists2} $V$ entorno abierto de $p\in S$.
        \item \label{def:sup_exists3} $X: U \to \bb{R}^3$ diferenciable
    \end{enumerate}
    tales que
    \begin{enumerate}[label=\arabic*)]
        \item \label{def:sup_cond1} $X(U)=V$
        \item \label{def:sup_cond2} $X: U \to V$ es un homeomorfismo
        \item \label{def:sup_cond3} $(dX)_q: \bb{R}^3\to \bb{R}^3$ es un monomorfismo $\forall q\in U$
    \end{enumerate}
\end{definicion}

\begin{observacion}
    La condición \ref{def:sup_cond3} es equivalente a decir que 
    \begin{gather*}
        \{X_u(q), X_v(q)\} \text{ son linealmente independientes} \quad \forall q\in U
    \end{gather*}
    que a su vez es equivalente a que
    \begin{gather*}
        X_u(q) \wedge X_v(q) \neq 0 \quad \forall q \in U
    \end{gather*}
\end{observacion}

\begin{notacion}
    Al plano tangente de un punto $p\in S$ lo notaremos por $T_pS$.
\end{notacion}

\section{Teorema de Separación de Jordan-Brower}

Para seguir un procedimiento análogo al estudio de Curvas del tema anterior nos tocaría ahora hacer una clasificación de superficies. Sin embargo no tenemos herramientas para este fin (este estudio se realiza en Topología II para las superficies compactas). Nos podemos preguntar por qué no tenemos herramientas. Esto se debe a que las topologías posibles para las superficies son infinitas (basta pensar en el $n$-toro) mientras que con curvas solo había 2 (recta y curva de Jordan).\\

Lo siguiente que hicimos tras el \hyperref[teo:clasif]{Teorema de Clasificación} fue el \hyperref[teo:curva_jordan]{Teorema de Jordan} para curvas para el cual necesitábamos dos ingredientes (transversalidad y \hyperref[teo:sard]{Teorema de Sard}). Veamos la definición de transversalidad para superficies.

\begin{definicion}
    Sea $S$ una superficie y $R$ una recta. Supongamos que $\exists p\in R\cap S$. 
    \begin{itemize}
        \item Diremos que $R$ es \textbf{tangente} a $S$ en $p$ cuando la variedad de dirección de la recta está en el plano tangente, es decir
        \begin{gather*}
            \vec{R}\subset T_pS
        \end{gather*}
        y escribiremos $R\ \|_p S$.

        \item Diremos que $R$ es \textbf{transversa} a $S$ en $p$ si y solo si 
        \begin{gather*}
            \vec{R} \not\subset T_pS
        \end{gather*}
        y escribiremos $R \transv_p S$.
    \end{itemize}
\end{definicion}

\begin{observacion}
    Las propiedades de la transversalidad son análogas al caso de las curvas simples.
\end{observacion}

Uniendo este concepto con un Teorema de Sard para dimensión $2$ llegamos al siguiente teorema:

\begin{teo}[Teorema de Separación de Jordan-Brower]\label{teo:sep_jordan_brower}
    Sea $S\subset \bb{R}^3$ una superficie conexa y cerrada en $\bb{R}^3$, entonces $\bb{R}^3 \setminus S$ tiene exactamente 2 componentes conexas. Si además $S$ es compacta, solo una de ellas es acotada (se llama el \textbf{dominio interior} determinado por $S$).
\end{teo}

La demostración de este teorema es análoga a la de \hyperref[teo:curva_jordan]{Teorema de la Curva de Jordan} pero para el caso de superficies.

\section{Orientabilidad}

En Curvas y Superficies estudiamos qué significa para una superficie dada $S\subset \bb{R}^3$ ser orientable. Lo recordamos a continuación

\begin{definicion}
    Sea $S\subset \bb{R}^3$ una superficie, diremos que $S$ es \textbf{orientable} si $\exists N :s \to \bb{R}^3$ \underline{diferenciable} tal que $\forall p\in S$ se tiene que 
    \begin{gather*}
        N(p)=N_p \perp T_pS \qquad \text{ con } \qquad |N_p| = 1
    \end{gather*}
    dicha $N$ se llama \textbf{aplicación de Gauss}.
\end{definicion}

\begin{observacion}
    Estudiamos ya que toda superficie es localmente orientable. Recordamos la demostración, ya que nos será de utilidad en lo que sigue
    \begin{proof}
        Consideramos $X: U \to S$ una parametrización y definimos $N^X : U \to \bb{R}^3$ dada por
        \begin{gather*}
            N^X = \frac{X_u \wedge X_v}{|X_u \wedge X_v|}
        \end{gather*}
        y es claro que $|N^X| = 1$. Además, es también claro por la definición de producto escalar que 
        \begin{gather*}
            N^X \perp \{X_u, X_v\} \Rightarrow X^X \perp T_XS
        \end{gather*}
        Sea $N:X(U) \to \bb{R}^3$ definimos
        \begin{gather*}
            N = N^X \circ X^{-1}
        \end{gather*}
        y ya tenemos la aplicación que buscábamos, ya que $N$ es claramente diferenciable por ser composición de diferenciables.
    \end{proof}
    Estudiamos también que no toda superficie es globalmente orientable (recordamos el ejemplo de la Cinta de Möbius, ya estudiado en la asignatura de Curvas y Superficies).
\end{observacion}

\subsection{Teorema de Brouwer-Samelson}

Vamos a ver ahora un resultado que nos garantizará la orientabilidad de superficies compactas (no necesariamente conexas).

\begin{teo}[Teorema de Brouwer-Samelson]\label{teo:brouwer_samelson}
    Toda superficie compacta es orientable.

    \begin{proof}
        Podemos suponer sin pérdida de generalidad que $S$ es conexa, ya que en caso contrario podemos trabajar en cada componente conexa suya.
        Vamos a tomar $p\in S$ y consideramos $T_pS$ que es un plano de $\bb{R}^3$. Sea $n_p$ a un vector perpendicular al tangente y con norma 1, es decir,
        \begin{gather*}
            n_p \perp T_pS \quad \text{ y } \quad |n_p|=1 
        \end{gather*}
        (Sabemos que hay otro con las mismas propiedades, $-n_p$). Sea $R$ la recta que pasa por $p$ en dirección $n_p$ (o $-n_p$). Sabemos que 
        \begin{gather*}
            p\in S\cap R \Rightarrow R \transv_p S
        \end{gather*}
        Por las propiedades de transversalidad sabemos que $\exists \veps>0$ de forma que 
        \begin{gather*}
            \{p+tn_p \ : \ |t|<\veps\}\cap S = \{p\}
        \end{gather*}
        Quitando el punto $p$ podemos separar este conjunto en 
        \begin{gather*}
            \{p+tn_p \ : \ 0<t<\veps\} \quad \text{ y } \quad \{p+tn_p \ : \ -\veps<t<0\}
        \end{gather*}
        Que son ambos conexos y no tocan a $S$. Por tanto tenemos 2 posibilidades:
        \begin{enumerate}
            \item $\{p+tn_p \ : \ 0<t<\veps\} \subset \Omega \quad$ (y así $\{p+tn_p \ : \ -\veps<t<0\}\subset \bb{R}^3 \setminus \overline{\Omega}$). En este caso diremos que $n_p$ ``apunta al interior'' de $S$, es decir, $\Omega$.
            \item $\{p+tn_p \ : \ 0<t<\veps\} \subset \bb{R}^3 \setminus \overline{\Omega} \quad$ (y así $\{p+tn_p \ : \ -\veps<t<0\}\subset \Omega$). En este caso diremos que $n_p$ ``apunta al exterior'' de $S$.
        \end{enumerate}
        Si $n_p$ apunta al interior tendremos que $-n_p$ apunta al exterior y viceversa.\\

        Definimos ahora $N:S\to \bb{R}^3$ como 
        \begin{gather*}
            % \text{ para cada } p\in S \text{ se tiene } N(p)=N_p =
            \left[
            \begin{array}{c}
                \text{único vector normal a $T_pX$}\\
                \text{que es unitario y apunta al}\\
                \text{interior de $S$}
            \end{array}
            \right]
            \qquad \forall p\in S
        \end{gather*}
        queremos ver ahora que $N$ es diferenciable. Para ello tomamos un $p_0\in S$ y $X:U\subset \bb{R}^2 \to S\subset \bb{R}^3$ una parametrización con
        \begin{gather*}
            p_0\in X(U) \qquad p_0=X(u_0,v_0) \qquad (u_0,v_0)\in U
        \end{gather*}
        Definimos ahora una nueva aplicación $G:U\times \bb{R} \to \bb{R}^3$ dada por\footnote{nos debería recordar a la aplicación exponencial que se definía para el estudio de los entornos tubulares en curvas.}
        \begin{gather*}
            G(u,v,t) = X(u,v) + tN^X(u,v)
        \end{gather*}
        Es claro que $G$ es diferenciable por serlo $X$ y $N^X$. Vamos a estudiar sus derivadas parciales:
        \begin{align*}
            \frac{\partial G}{\partial u} (u,v,t) &= G_u(u,v,t) = X_u(u,v) + t N^X_u(u,v)\\[1em]
            \frac{\partial G}{\partial v} (u,v,t) &= G_v(u,v,t) = X_v(u,v) + t N^X_v(u,v)\\[1em]
            \frac{\partial G}{\partial t} (u,v,t) &= G_t(u,v,t) = N^X(u,v)
        \end{align*}
        Si estudiamos su comportamiento en $t=0$ tenemos
        \begin{gather*}
            G_u(u,v,0) = X_u(u,v)\qquad
            G_v(u,v,0) = X_v(u,v)\qquad
            G_t(u,v,0) = N^X(u,v)
        \end{gather*}
        Si observamos estos vectores vemos que son linealmente independientes ya que $\{X_u, X_v\}$ forman una base de $T_{X(u,v)}S$ y al añadirle el normal (que es claramente linealmente independiente al tangente) tendremos que forman una base de $\bb{R}^3$. Tendríamos la siguiente transformación
        \begin{align*}
            (dG)_{(u,v,0)} &: \bb{R}^3 \to \bb{R}^3\\
            &(1,0,0) \mapsto X_u(u,v)\\
            &(0,1,0) \mapsto X_v(u,v)\\
            &(0,0,1) \mapsto N^X(u,v)
        \end{align*}
        que lleva una base de $\bb{R}^3$ en otra. Por tanto $(dG)_{(u,v,0)}$ es regular. En particular, $(dG)_{(u_0,v_0,0)}$ es regular y aplicando el Teorema de la Función Inversa tenemos que existe un entorno de $(u_0, v_0, 0)$ en $\bb{R}^3$, que podemos elegir de la forma
        \begin{gather*}
            W \times (-\veps, \veps)
        \end{gather*}
        donde $\veps>0$ y $W$ es un entorno abierto de $(u_0, v_0)$ conexo\footnote{Si no fuese conexo nos quedaríamos con la c.c. que tuviese al punto $(u_0, v_0)$} ($W\subset U$). Existe también un entorno abierto de $G(u_0, v_0, 0)=p_0$ en $\bb{R}^3$ tales que $G(W\times(-\veps, \veps))=\tilde{W}$ de forma que 
        \begin{gather*}
            G_{|_{W\times(-\veps, \veps)}} : W\times (-\veps, \veps) \to \tilde{W}
        \end{gather*}
        es un difeomorfismo. Tenemos que 
        \begin{gather*}
            W\times\{0\} \overset{G_{|_{W\times(-\veps, \veps)}}}{\longmapsto} S\cap \tilde{W} \qquad \text{ entorno de $p_0$ en $S$}
        \end{gather*}
        De esta forma tenemos que 
        \begin{align*}
            \tilde{W}-S &= G(W\times(-\veps, \veps)) - G(W\times\{0\})\\
            &=
            G(W\times(0,\veps)) \dot{\cap} G(W\times(-\veps,0))
        \end{align*}
        Hacemos ahora una distinción de casos:
        \begin{enumerate}
            \item[i)] Si $G(W\times (0,\veps))\subset \Omega$, entonces tendríamos
            \begin{gather*}
                X(u,v) + tN^X(u,v) \qquad (u,v)\in W \quad 0<t<\veps
            \end{gather*}
            Por lo que $N^X$ apunta al interior en todo $W$, de forma que $N^X\circ X = N_{|_W}$.
            
            \item[ii)] Si $G(W\times (0,\veps))\subset \bb{R}^3 \setminus \overline{\Omega}$, entonces $N^X$ apunta al exterior por lo que 
            \begin{gather*}
                -N^X \circ X^{-1} = N
            \end{gather*}
        \end{enumerate}
        en cualquiera de los casos tendríamos que es diferenciable, lo que concluye la demostración.
    \end{proof}
\end{teo}

